\section{Proposed Method}
\label{sec:method}

The method grew in three stages, each answering a failure of the previous one.
\Cref{sec:method:wrapper} wraps a frozen model in cross-image conformal calibration (broad study).
\Cref{sec:method:anchors} adds sparse test-time anchors and a distance-aware model of their information (broad study).
\Cref{sec:method:dual} presents the main method, which separates the anchors' two roles.
\Cref{tab:notions} in \cref{app:notions} summarizes which method gives which guarantee.

\subsection{A conformal wrapper for frozen depth models}
\label{sec:method:wrapper}
All intervals have the form $\yhat\,e^{\pm\lambda s(x)}$ (shifted by $m$ where anchors are available), with a \emph{scale map} $s$ that encodes how difficult each pixel is, and a threshold $\lambda$ set by calibration.

\paragraph{Scale maps and calibrators.}
We compare nine scale maps: constant error in metres, error proportional to $\yhat$, constant relative error (the \emph{plain log score}, $s\equiv1$), flip disagreement, ensemble spread, gradient-based uncertainty~\cite{hornauer2022gradient}, the model's own confidence map, a \emph{learned} score (gradient-boosted trees fitted on held-out \emph{head-train} images), and a \emph{factorised} score $a(\text{image})\times b(\text{pixel})$ that separates ``how wrong is this image'' from ``where in it''.
With images as units, the calibrators (\cref{tab:calibrators} in \cref{app:broad}) are: K1, conformal risk control on the fraction of an image's pixels left uncovered (marginal coverage over images); RCPS~\cite{bates2021rcps}, its high-probability version; V4, split conformal on each image's $(1-\beta)$-quantile of scores (the image event); K2, every pixel of a fixed patch covered at once; K3, Mondrian per predicted-depth bin or for edge versus smooth pixels; and K0, a naive baseline that treats pixels as exchangeable.
Scene-level versions (K1-scene, V4-scene, K2-scene) weight scenes equally and are valid when scenes, not images, are exchangeable; they need at least $1/\alpha-1$ calibration scenes.

\subsection{Sparse test-time anchors}
\label{sec:method:anchors}
Each test image comes with $k$ anchors, and calibration images receive anchors by the same rule.
In the broad study all $k$ measurements play both roles.
\begin{itemize}[leftmargin=*,itemsep=1pt,topsep=2pt]
\item \mRescale shifts the prediction by the median anchor residual, $m=\med_{p\in\cA}\dtil(p)$, then calibrates across scenes (K1-scene). It gives a marginal guarantee.
\item \mHybrid additionally divides each pixel's score by the image's anchor-estimated spread (leave-one-out median absolute deviation plus a floor $\tau$ from head-train). It gives a marginal guarantee.
\item \mLocal calibrates inside the image only: half of the anchors fix the shift, the other half set the width by split conformal~\cite{mao2024spatial}.
If the anchors are a uniform sample of pixels, it gives a per-image guarantee with exact value $r/(k/2+1)$, $r=\lceil(k/2+1)(1-\alpha)\rceil$, and no calibration set is needed.
It is infinite for $k<18$ at $\alpha=0.1$.
\item \mBoth is the pointwise hull of \mHybrid and \mLocal; since both statements are monotone in the interval, it inherits both guarantees~\cite{gasparin2024merging}.
\end{itemize}

\paragraph{Distance-aware anchors (\mDist).}
Under structured layouts, coverage depends strongly on the distance to the nearest anchor (\cref{sec:res:broad}).
\mDist models this.
The local shift $m(p)$ is one of four estimates from the anchor residuals: the image median, the median of the 16 nearest anchors, or a Gaussian or bilateral (also weighting similarity in $\log\yhat$) kernel mean.
The local spread $\sigma(p)$ is the median absolute deviation of the 16 nearest anchor residuals, and the scale is
\begin{equation}
s(p)=\big(\sigma(p)+\tau\big)\,g\big(\delta(p)\big),
\label{eq:scale}
\end{equation}
with a floor $\tau$ and a non-decreasing $g$ fitted by isotonic regression of normalized error against distance.
The variant and bandwidths are chosen per layout on head-train images (``auto''), never on calibration or test data.
One threshold per bin of $\delta$ is then set by CRC over calibration scenes, with the miscovered fraction of bin pixels as loss.
This is a Mondrian construction: the expected bin miscoverage of a new scene is at most $\alpha$ in every bin, and sparse bins are merged.
In the main method this calibrated version is called \mDistSplit, and the same anchor model supplies the centre of \mDual.

\paragraph{Layout mismatch.}
Cross-image calibration is valid only if calibration and test anchors follow the same layout.
If the deployed layout is one of a known set, calibrating under each and keeping the largest threshold (``@max'') is valid for each of them: for every layout, the @max threshold is at least that layout's own valid threshold.
A random mixture of layouts at calibration time is valid only on average over that mixture.

\subsection{Main method: dual-role anchors (\mDual)}
\label{sec:method:dual}

%% Pipeline figure (TikZ, single column, drawn at native width ~7.9 cm).
\begin{figure}[t]
\centering
\begin{tikzpicture}[
  font=\sffamily\scriptsize,
  box/.style={draw, rounded corners=2pt, align=center, text width=1.62cm, minimum height=8mm, inner sep=2pt, fill=white},
  frozen/.style={box, fill=gray!15},
  anch/.style={box, fill=blue!8},
  prob/.style={box, fill=orange!18},
  outbox/.style={box, text width=1.75cm, fill=white, very thick},
  opt/.style={box, text width=1.75cm, dashed, fill=green!6},
  arr/.style={-{Stealth[length=1.6mm]}, semithick},
  darr/.style={-{Stealth[length=1.6mm]}, semithick, dashed}]
  % column 1: inputs
  \node[box]    (rgb)    at (0.85, 2.15) {RGB image $x$};
  \node[box]    (sensor) at (0.85, 0.45) {sparse depth,\\budget $k$};
  % column 2: model and split
  \node[frozen] (f)      at (2.85, 2.15) {frozen model $f$\\[-1pt]{\tiny(or any probe-indep.\ centre)}};
  \node[anch]   (S)      at (2.85, 0.95) {$k{-}n$ structured\\anchors $\cS$};
  \node[prob]   (Q)      at (2.85,-0.15) {$n$ uniform\\probes $\cQ$};
  % column 3: centre/scale and threshold
  \node[anch]   (ms)     at (4.85, 2.15) {centre $m$, scale $s_\gamma$\\[-1pt]{\tiny from $\cS$ only}};
  \node[prob]   (q)      at (4.85,-0.15) {threshold\\$\qhat_i=R_{(r)}$};
  % column 4: output and optional stages
  \node[outbox] (int)    at (6.95, 2.15) {$\yhat\,e^{\,m\pm\qhat_i s_\gamma}$\\[-1pt]\textbf{\mDual}};
  \node[opt]    (hull)   at (6.95, 0.95) {$\cup$ \mDistSplit\\[-1pt]\textbf{\mDualBoth}};
  \node[opt]    (env)    at (6.95,-0.15) {sensor-error\\envelope};
  % arrows
  \draw[arr] (rgb) -- (f);
  \draw[arr] (f) -- node[above, font=\tiny\sffamily]{$\yhat$} (ms);
  \draw[arr] (sensor.east) -- ++(0.12,0) |- (S.west);
  \draw[arr] (sensor.east) -- ++(0.12,0) |- (Q.west);
  \draw[arr] (S.east) -- (ms.south west);
  \draw[arr] (Q) -- (q);
  \draw[arr] (ms) -- (q);
  \draw[arr] (ms) -- (int);
  \draw[arr] (q.east) -- ++(0.12,0) |- ([yshift=-2mm]int.west);
  \draw[darr] (int) -- (hull);
  \draw[darr] (hull) -- (env);
  \node[font=\tiny\sffamily, text=orange!70!black, anchor=north] at (2.85,-0.62) {never seen by $f$, $m$, $s$};
\end{tikzpicture}
\caption{\textbf{Dual-role calibration.} Structured anchors $\cS$ build a centre $m$ and scale $s_\gamma$ around a frozen prediction $\yhat$ (or any probe-independent centre); the $n$ uniform probes $\cQ$ only score it and set one threshold per image (\cref{thm:dual}). Dashed: optional stages (\ifenvelopemain\cref{thm:both,prop:env}\else\cref{thm:both}; Supp.~D\fi).}
\label{fig:pipeline}
\end{figure}


\paragraph{Architecture (\cref{fig:pipeline}).}
(1)~A frozen monocular model produces $\yhat$.
(2)~The sensor budget is split into $k-n$ anchors $\cS$ that follow the deployed layout and $n$ uniformly random probes $\cQ$.
(3)~The anchors alone build a local centre $m(p)$ and scale $s(p)$ around $\yhat$ with the \mDist model of \cref{eq:scale}, fitted on $\cS$ only.
(4)~The probes, unseen by stages 1--3, set one conformal threshold per image.
Stage~1 can be any predictor that does not see the probes, \eg a depth-completion network fed RGB and $\cS$; stage~3 is then refit around its output.
\mDual is therefore a calibration layer, not a depth estimator.

\paragraph{Scale shrinkage.}
When the anchors set the shift wrong everywhere, as a single cluster can, the residual is nearly flat.
The threshold is then set by the probes with the smallest $s$, and widths elsewhere grow by up to $\max s/\min s$.
We therefore shrink the scale toward its per-image median $M_i=\med_{p\in\cP_i}s(p)$:
\begin{equation}
s_\gamma(p)=M_i\big(s(p)/M_i\big)^{\gamma},\qquad\gamma=0.5.
\label{eq:shrink}
\end{equation}

\paragraph{Per-image threshold.}
For probes $\cQ=\{p_1,\dots,p_n\}$ with reported residuals $\dtil(p_j)$,
\begin{equation}
R_j=\frac{|\dtil(p_j)-m(p_j)|}{s_\gamma(p_j)},\qquad\qhat_i=R_{(r)},\qquad r=\lceil(n+1)(1-\alpha)\rceil,
\label{eq:threshold}
\end{equation}
where $R_{(r)}$ is the $r$-th smallest score ($\qhat_i=\infty$ if $r>n$).
The interval at every evaluation pixel is
\begin{equation}
C(p)=\Big[\yhat(p)\,e^{\,m(p)-\qhat_i s_\gamma(p)},\;\yhat(p)\,e^{\,m(p)+\qhat_i s_\gamma(p)}\Big].
\label{eq:interval}
\end{equation}

\begin{theorem}[Dual-role validity]
\label{thm:dual}
Fix an image with valid pixels $\cP$ and a structured anchor set $\cS\subset\cP$ chosen by any rule, which may depend on $x$, $\yhat$ and the anchor values.
Let $m$, $s>0$ and $\gamma$ be determined by $(x,\yhat,\cS,\dtil|_\cS)$.
Let $v$ be the log residual of the quantity the probes report ($v=d$ for exact depth, $v=\dtil$ for a sensor).
Draw $\cQ$ uniformly without replacement from $\cU=\cP\setminus\cS$, compute $\qhat$ by \cref{eq:threshold} with $v$, and let $J$ be uniform on $\cU\setminus\cQ$.
Then
\begin{equation*}
\P\big(\yhat(J)e^{v(J)}\in C(J)\,\big|\,x,\cS,\dtil|_\cS\big)\;\ge\;\frac{r}{n+1}\;\ge\;1-\alpha,
\end{equation*}
with equality on the left when the scores have no ties and $r\le n$.
Equivalently, $\E[\Cov_i\mid x,\cS]\ge r/(n+1)$, the expectation being over the probe draw.
\end{theorem}

\noindent\emph{Proof idea.}
Given the anchors, every pixel of $\cU$ carries a fixed score.
The probes together with $J$ form a uniformly random $(n+1)$-subset of $\cU$, so $J$'s score is among the $r$ smallest of the $n+1$ scores, and hence covered, with probability at least $r/(n+1)$.
The full proof is in \cref{app:proof-dual}. \hfill$\square$

The theorem needs no calibration images and uses all $n$ probes for the threshold; $n\ge9$ suffices at $\alpha=0.1$.
By contrast, \mLocalSplit, which applies \mLocal to the $n$ probes, spends half of them on a shift.
\emph{Scope:} the guarantee holds in expectation over the probe draw, not per region.
Probes must be uniform over the pixels it concerns.
Every choice, $\gamma$ included, must be fixed without looking at probe values.
And it covers $v$, the quantity the probes measure.

\paragraph{Probe budget.}
\Cref{thm:dual} fixes the mean of $\Cov_i$; the number of probes fixes its spread.

\begin{proposition}[Per-image coverage law]
\label{prop:budget}
Under \cref{thm:dual} with no ties and $r\le n$, let $N=|\cU|$.
The number of covered evaluation pixels, $(N-n)\Cov_i$, follows exactly the beta-binomial law $\BetaBin(N-n,\,r,\,n+1-r)$.
As $N\to\infty$, $\Cov_i\to\BetaD(r,\,n+1-r)$ in distribution.
\end{proposition}

This is the finite-population form of the Beta law of split-conformal coverage conditional on the calibration set~\cite{vovk2012conditional}; the proof is in \cref{app:proof-budget}.
It turns the probe count into a design parameter (\cref{tab:budget}).
The share of images below coverage $c$ in one draw, $\P(\BetaD(r,n{+}1{-}r)<c)$, is known before any data are seen.
At $\alpha=0.1$ it falls slowly with $n$ (\NumRuleShortHundred\% below 0.85 at $n=100$).
The nominal level is the stronger lever: $r=n=20$ (the \mDualRmax arm, nominal \NumRuleNomTwentyAll) leaves only \NumRuleShortTwentyAll\% of images below 0.85, at a width cost the confirmatory run measures (\pending{measured, $r=n$ arm}).

\begin{table}[t]
\centering\small
\caption{\textbf{Probe budget at $\alpha=0.1$, one probe draw}, from the Beta law of \cref{prop:budget} alone (no data).}
\label{tab:budget}
\begin{tabular}{@{}rrccc@{}}
\toprule
$n$ & $r$ & $r/(n+1)$ & 10th percentile of $\Cov_i$ & $\P(\Cov_i<0.85)$, \% \\
\midrule
10 & 10 & 0.909 & \NumBetaPTenTen   & \NumBetaShortOneTen \\
20 & 19 & 0.905 & \NumBetaPTenTwenty & \NumBetaShortOneTwenty \\
40 & 37 & 0.902 & \NumBetaPTenForty & \NumBetaShortOneForty \\
\bottomrule
\end{tabular}
\end{table}

\paragraph{A per-distance-bin guarantee (\mDualBoth).}
\mDual's coverage, in practice, decreases with the distance to the nearest anchor; \cref{thm:dual} allows this because it is an average over the image.
\mDistSplit (\cref{sec:method:anchors}, fitted on all $k$ measurements) has the opposite profile: a per-bin guarantee averaged over scenes, no per-image guarantee, and a need for calibration scenes with dense depth from the same source as the measurements.
\mDualBoth is the pointwise hull
$C^{\text{both}}(p)=\big[\min(\ell^{\text{dual}},\ell^{\text{dist}}),\,\max(u^{\text{dual}},u^{\text{dist}})\big]$.

\begin{theorem}[Hull validity]
\label{thm:both}
\mDualBoth has the per-image guarantee of \cref{thm:dual} and the per-bin guarantee of \mDistSplit.
\end{theorem}

\noindent\emph{Proof idea.} The hull contains each interval pointwise and both statements are monotone in the interval~\cite{gasparin2024merging} (\cref{app:proof-both}). \hfill$\square$

\noindent The broad study's \mDistBoth, the hull of \mDistSplit with \mLocalSplit, is the baseline that \mDualBoth replaces.
Any hull costs width, most when its two centres disagree, as under clustered anchors.

\subsection{Sensor-error envelope}
\label{sec:env}
With sensor-valued probes, \cref{thm:dual} certifies the reading $\ytil$, not the depth $y$.
Let $e=\log\ytil-\log y$ and suppose a reference (here a laser scan) exists on calibration scenes.
We (a)~subtract a global bias $\hat\beta$, the median of $e$ on head-train pixels, from the probe values.
We (b)~give each sensor-confidence class $c$ (ARKit: $0,1,2$) the band $[-a_c,b_c]=[Q_c(\lambda/2),Q_c(1-\lambda/2)]$ of the corrected head-train error, clipped to contain 0.
One conformal risk control over the $U$ calibration scenes sets $\hat\lambda$: the largest value on a fixed grid with $\frac{U}{U+1}\hat L(\lambda)+\frac{1}{U+1}\le\alpha_e$, where $\hat L$ is the mean fraction of a scene's pixels outside their band.
With $[\ell,u]$ the \mDual interval for the corrected reading at level $1-\alpha_s$, the enveloped interval (\mDualEnv) is
\begin{equation}
C^{\text{env}}(p)=\big[\ell(p)\,e^{-b_{c(p)}},\;u(p)\,e^{a_{c(p)}}\big].
\label{eq:envelope}
\end{equation}

\begin{proposition}[Envelope]
\label{prop:env}
If the sensor reading and its confidence are defined at every evaluation pixel, the miscoverage of $y$ by $C^{\text{env}}$, averaged over test scenes and probe draws, is at most $\alpha_s+\alpha_e$.
\end{proposition}

\noindent\emph{Proof idea.} $y$ can be missed only if the corrected reading falls outside $[\ell,u]$ (at most $\alpha_s$ per image, \cref{thm:dual}) or the corrected error falls outside its band (at most $\alpha_e$ in expectation, by risk control); a union bound concludes (\cref{app:proof-env}). \hfill$\square$

\noindent The per-image guarantee on $y$ is not retained, since the envelope part holds on average over scenes.
We fixed $\alpha_s=\alpha_e=0.05$ before the run.
One risk control over a single class-adaptive parameter replaces per-class thresholds, which would all be infinite with 37 calibration scenes (\cref{app:proof-env}); this design was amended before any envelope output existed.

